\chapter*{Abstract}
\addcontentsline{toc}{chapter}{Abstract}
\markboth{Abstract}{Abstract}

This Bachelor's Thesis presents an experimental comparison between two strategies for analyzing electrocardiograms represented as images under data scarcity. The central question is direct: given the same budget of \(K\) labeled examples, when is it worth training a task-specific vision model (CNN) and when is it preferable to leverage the multimodal knowledge of a vision-language model (VLM) through in-context learning? The motivating clinical case is suspected Brugada syndrome, a hereditary sodium channel disorder associated with sudden cardiac death.

The work is strictly framed as clinical triage support. In routine practice, when an ECG suggests a suspicious Brugada pattern, a clinical circuit is initiated that includes expert ECG reading, assessment of symptoms and family history, possible repositioning of leads V1 and V2 in higher intercostal spaces (second or third instead of the standard fourth), and, in selected patients with persistent suspicion and a non-diagnostic baseline ECG, a pharmacological challenge with sodium channel blockers (ajmaline, flecainide, procainamide, or pilsicainide) in a monitored setting with advanced cardiac life support capabilities. This test can unmask a Type 1 pattern but must be interpreted within the clinical context and after excluding phenocopies. The model developed here does not diagnose or rule out patients: it prioritizes traces that deserve early expert review.

The work integrates two domains. The first is a controlled QRS/ST morphology simulator that renders V1, V2 and V3 images with labels from observable rules. The second is the Brugada HUCA dataset, processed from real clinical recordings with 317 included patients, 951 images and a binary label of confirmed Brugada case versus control. Evaluation is patient-level leave-one-out cross-validation with reduced budgets \(K=\{2,4,8,16,32\}\) and three seeds per budget.

The CNN branch serves a deliberate purpose: to empirically demonstrate that the classical supervised approach does not work with so few data, even with domain adaptation techniques. ResNet18 learns clearly on the simulator, reaching balanced accuracy 0.848 at \(K=32\), sensitivity 0.883 and specificity 0.812. On Brugada HUCA, however, performance drops close to chance: the best baseline balanced accuracy is 0.516 at \(K=16\). Domain adaptation (SSL, CORAL, MMD, DANN) improves the best point to 0.549 with MMD at \(K=32\), an improvement that remains insufficient for a robust system.

This negative evidence is a central contribution of the thesis. It is not a project failure but the result that grounds the main hypothesis: if training a task-specific visual model is not enough with few data, the prior multimodal knowledge of a VLM could compensate for that scarcity. On that basis, the VLM/ICL branch is developed as the main comparator: Gemma 4 and MedGemma receive the same images and the same \(K\) budgets, but examples are included as in-context demonstrations without updating weights. The document includes the full protocol, prompts, JSON validation, controls and prepared result tables; only the numerical VLM cells are marked as \emph{to be run} so that no inference results are invented.
